\section{Implementing the Upward, Horizontal, and Downward Passes}
\label{sec:translation-passes}

With our octree generated, the near-field and far-field interaction lists populated, and the dense translation operators fully precomputed, we can now begin the hierarchical traversal of the octree to compute the far-field interactions. 
This traversal happens in every iteration of iterative solver used to solve the system.
There are three distinct, sequential passes in the traversal: the upward pass, the horizontal pass, and the downward pass.

\subsection{Upward Pass: Aggregating Far-Field Signatures} \label{upward-pass}
The traversal begins with the upward pass. It evaluates the initial acoustic influence of the boundary sources contained in the active solution vector (in our case, the $\phi$ estimated at the previous iteration) in all the leaf nodes. For each leaf node, the \ac{P2M} operator is multiplied by the specific potential and flux (in our case the normal velocity from our Right-hand side) values assigned to the local boundary elements. This operation yields the multipole expansion coefficients for that leaf node.

The algorithm then traverses up the octree, level by level, from the leaf nodes to the root node. For each parent node, its generalized multipole expansion coefficients are computed by translating the multipole expansions of its child nodes to its own geometric center. Because the \ac{M2M} operator applies the addition theorem linearly, the parent's total multipole expansion is evaluated as the algebraic sum of the shifted multipole expansions of its children.

\newglossaryentry{Mparent}{
  name={$\mathbf{M}_{\mathrm{parent}}$},
  description={Multipole expansion coefficient vector of the parent node}
}
\newglossaryentry{Tm2m}{
    name={$\mathbf{T_{M2M}}$},
    description={Dense translation matrix for Multipole-to-Multipole operator}
}
\newglossaryentry{Mchild}{
    name={$\mathbf{M}_{\mathrm{child}}$},
    description={Multipole expansion coefficient vector of a child node}
}
\newglossaryentry{cn}{
    name={$c_n$},
    description={Number of child nodes of a parent node}
}

\begin{equation}
    \mathbf{M}_{parent} = \sum_{c=1}^{c_n} \mathbf{T}_{M2M}^{(c)} \cdot \mathbf{M}_{child}^{(c)}
\end{equation}
where $c_n$ is the number of child nodes.
\subsection{Horizontal Pass: Translating Multipole Expansions to Local Expansions} \label{horizontal-pass}

The horizontal pass serves as the critical bridge of the \ac{FMM} algorithm, converting the outgoing scattered radiation represented by the multipole expansions into incoming incident fields represented by the local expansions.
For all valid nodes with depth $d\geq 2$, the horizontal pass iterates through the far-field interaction list of the node. For each target-source pairing within the list, the source node's multipole expansion is multiplied by the specific precomputed \ac{M2L} translation operator associated with the relative displacement vector separating the source and target nodes. 

\newglossaryentry{Ltarget}{
    name={$L_{target}$},
    description={Local expansion coefficient vector of the target node}
}
\newglossaryentry{Tm2l}{
    name={$\mathbf{T}_{M2L}$},
    description={Dense translation matrix for Multipole-to-Local operator}
}
\newglossaryentry{Msource}{
    name={$\mathbf{M}_{source}$},
    description={Multipole expansion coefficient vector of the source node}
}


The resulting localized coefficients, which represent the acoustic contribution arriving from distant source nodes are algebraically accumulated into the target node's local coefficient array:
\begin{equation}
    \mathbf{L}_{target} \mathrel{+}= \mathbf{T}_{M2L} \cdot \mathbf{M}_{source}
\end{equation}
Due to us enforcing a uniform truncation order $P$, the dimensions of all the vectors and matrices involved line up, enabling the use of vectorized matrix-vector multiplication for a significant boost in computational efficiency.

\subsection{Downward Pass: Distributing Local Expansions to Leaf Nodes} \label{downward-pass}
The downward pass cascades the accumulated incident local expansion from the higher levels of the octree down to the leaf nodes, where the final evaluation of the far-field contribution to the acoustic field at the boundary elements takes place. 
 Starting from the nodes on level $d=2$, and traversing downward, the parent node's aggregated local coefficients are geometrically shifted to the centers of its child nodes using the precomputed \ac{L2L} translation operator. These shifted coefficients represent the macroscopic background acoustic field, are added to the existing local coefficient the child node accumulated during its horizontal pass interactions. This recursive cascade ensures that when the traversal reaches the leaf nodes, they have a complete representation of the far-field influence from all distant sources.

 The final step of the downward pass occurs at the leaf nodes, where the \ac{L2P}
 operator transforms the leaf node's local coefficients into the physical acoustic potential and flux (normal velocity) at the specific Cartesian coordinates of the target boundary elements.

 Operating in local spherical coordinates relative to the leaf center, the acoustic potential at a target point can be obtained by a straightforward double summation like so:
\begin{equation}
    \text{potential} = \sum_{n=0}^{P} \sum_{m=-n}^{n} L_{coeffs}\left[n^2 + n + m \right] \cdot j_n(kr) \cdot Y_n^m(\theta, \phi)
\end{equation}

The normal derivative calculation follows precisely the inverse of the \ac{P2M} logic.
The radial derivative uses the exact stable recurrence $j_n' = k \frac{n j_{n-1}-j_{n+1}}{2n +1}$, and the angular gradients use the same quantum mechanical ladder operators. The final physical normal derivative is generated by the dot product of this resolved spherical gradient vector with the target element's distinct surface normal vector.

Similar to our \ac{P2M} implementation, if the physical target point is too close to the leaf node center ($r < 10^{-8}$), we fall back to a \ac{FDM} approach to avoid the singularity. The final assembled physical basis vector output merges both the evaluated potential and the normal derivative, accurately scaled by the prefactor $-ik$ associated with the 3D Helmholtz equation.


\section{Building the Near-Field Matrix} \label{sec:near-field}

While the hierarchical translation passes of the \ac{FMM} efficiently approximate the far-field interactions, the near-field interactions must be computed with high precision to prevent the localized degradation of the solution. 
The near-field domain for any target node consist of the boundary elements in the node box itself, but also all the boundary elements in its geometrical neighbor boxes, and any coarser leaf boxes that are geometrically adjacent to the target node's parent node.

The acoustic interactions between the boundary elements in the near-field domain are computed directly with the \acf{P2P} mathematical integration. 
This integration logic is exactly the same as the one used to generate the dense system matrix for the \ac{BEM}.